\documentclass[11pt]{article}

\usepackage[margin=0.75in]{geometry}
\usepackage{booktabs}
\usepackage{amsmath}
\usepackage{multirow}
\usepackage{hyperref}

\title{Joint Detection-Loss Training: MMD vs.\ No-MMD, Cross-Domain Evaluation}
\author{}
\date{\today}

\begin{document}
\maketitle

\section{Setup}

Unlike every other comparison in this study, these checkpoints were trained with
\texttt{scripts/joint\_train.py}: both domains receive their own detection loss every
step (\texttt{mmd\_cfg.joint\_detection\_loss=True}), not just the target domain, with
the source domain forward-only as in the rest of this study. The MMD variant adds the
usual RBF-kernel alignment term (\texttt{mmd\_weight=0.8}, EMA bandwidth momentum 0.9,
\texttt{mmd\_target\_layer=10}, flattened features) on top of that; the no-MMD variant
is the joint detection loss alone.

Source = real\_small or syn\_small (both project's smallest subsets, 1600 train images
each), target = the other domain, run in both directions. YOLOv10n, imgsz 1920, batch 8,
100 max epochs, 2$\times$RTX~4090 via DDP (\texttt{--device}) per run. Every resulting
checkpoint was then evaluated on \emph{both} val sets -- real\_small val (548 images)
and syn\_large val (640 images, the val split syn\_small's own config borrows rather
than using its own 160-image val folder) -- with person and vehicle
reported separately, using a plain \texttt{model.val()} call (imgsz 1920) independent of
which domain that checkpoint was actually trained toward as its target.

\section{Results}

\begin{table}[h!]
    \centering
    \small
    \begin{tabular}{llrrrr}
        \toprule
        Model & Direction & Val set & person mAP50 & person mAP50-95 & vehicle mAP50 & vehicle mAP50-95 \\
        \midrule
        \multirow{4}{*}{MMD}
            & real$\to$syn & real & 0.574 & 0.263 & 0.854 & 0.602 \\
            & real$\to$syn & syn  & 0.604 & 0.294 & 0.611 & 0.423 \\
            & syn$\to$real & real & 0.530 & 0.240 & 0.831 & 0.582 \\
            & syn$\to$real & syn  & 0.614 & 0.292 & 0.612 & 0.424 \\
        \midrule
        \multirow{4}{*}{no-MMD}
            & real$\to$syn & real & 0.571 & 0.259 & 0.851 & 0.601 \\
            & real$\to$syn & syn  & 0.588 & 0.281 & 0.599 & 0.413 \\
            & syn$\to$real & real & 0.502 & 0.227 & 0.828 & 0.580 \\
            & syn$\to$real & syn  & 0.589 & 0.264 & 0.582 & 0.401 \\
        \bottomrule
    \end{tabular}
    \caption{Every joint-trained checkpoint evaluated on both val sets, person/vehicle
    separated. Columns are missing an overall/all-class average deliberately -- it is a
    weighted mix of both classes and not the quantity of interest here (as elsewhere in
    this study).}
\end{table}

\noindent MMD beats no-MMD on every single row -- both classes, both val sets, both
directions (8/8). The gap is largest on person in the syn$\to$real direction evaluated
on real val (0.530 vs.\ 0.502 mAP50, +0.028), and smallest -- effectively tied -- on
vehicle/real-val in the real$\to$syn direction (0.854 vs.\ 0.851). Since every checkpoint
here is evaluated on both val sets regardless of which domain was its actual adaptation
target, this table also gives a same-model, alternate-domain comparison for free.

\section{Compared to source-domain retention under the old (non-joint) regime}

Under the standard target-only-detection-loss adaptation used everywhere else in this
study, checkpoints score drastically worse on their source domain than on their
target domain (see \texttt{source\_domain\_retention.tex}: e.g.\ mAP50 dropping from the
$\sim$0.6--0.9 target-domain range to $\sim$0.02--0.35 on source). Joint training does
not show that collapse: every checkpoint above scores in a broadly similar,
reasonable range on \emph{both} val sets rather than one domain being near-zero. This is
the expected consequence of giving the source domain a real detection-loss gradient
every step instead of treating it as a frozen MMD reference -- but it has not yet been
checked whether this also holds at the project's standard scale (real\_large/syn\_large
or larger), only at this small-dataset, 100-epoch, 2-direction dev run.

\section{Open items}

\begin{itemize}
    \item Only the project's smallest subsets (real\_small/syn\_small) have been run
    through joint training so far -- no confirmation yet that MMD's edge over no-MMD, or
    the retained-both-domains pattern above, holds at real\_large/syn\_large/syn\_xlarge
    scale.
    \item No comparison yet against this study's main (non-joint) adaptation regimes at
    the same source/target pair -- i.e.\ whether joint training's target-domain numbers
    are competitive with, better, or worse than the target-only-detection-loss numbers
    reported elsewhere in this study for the same real\_small/syn\_small pair.
    \item Only one seed/run per (variant, direction) cell; no repeated-run variance
    estimate yet for how much of the MMD-over-no-MMD gap is noise vs.\ signal.
\end{itemize}

\end{document}
